\begin{figure}[t]
  \centering
  \begin{tikzpicture}[scale=0.91, transform shape,
    doc/.style={
      rectangle, rounded corners=3pt, draw=ThesisNeutral!65,
      fill=ThesisPaper, line width=0.65pt, text width=3.05cm,
      minimum height=2.25cm, align=left, inner sep=5pt,
      font=\scriptsize
    },
    field/.style={
      rectangle, rounded corners=3pt, draw=ThesisPrimaryDark,
      fill=ThesisPrimary!12, line width=0.7pt, text width=3.05cm,
      minimum height=1.25cm, align=left, inner sep=5pt,
      font=\scriptsize
    },
    map/.style={-{Stealth[length=1.7mm]}, line width=0.75pt,
      draw=ThesisNeutral, shorten <=2pt, shorten >=2pt},
    every node/.style={text=ThesisInk}
  ]
    \node[doc] (d1) at (-5.4,1.45) {%
      \textbf{12 mai 2018 · consultation}\\[3pt]
      Suspicion de lymphome; bilan diagnostique engagé.
    };
    \node[doc] (d2) at (-1.8,1.45) {%
      \textbf{22 mai 2018 · RCP}\\[3pt]
      Décision: début d'une \textcolor{ThesisPrimaryDark}{\textbf{chimiothérapie par protocole CHOEP-14}}, 8 cycles.
    };
    \node[doc] (d3) at (1.8,1.45) {%
      \textbf{5 octobre 2018 · TEP-TDM}\\[3pt]
      Réponse complète; score de Deauville~1.
    };
    \node[doc] (d4) at (5.4,1.45) {%
      \textbf{18 nov. 2033 · RCP}\\[3pt]
      Rappel de la RCP du 25 avr. 2032: après rechute, réintroduire 6 cycles de
      \textcolor{ThesisPrimaryDark}{\textbf{CHOEP-14}}.
    };

    \draw[line width=0.8pt, draw=ThesisNeutral!65]
      (-6.9,-0.35) -- (6.9,-0.35);
    \foreach \x/\date in {-5.4/12.05.18,-1.8/22.05.18,1.8/05.10.18,5.4/18.11.33} {
      \draw[line width=0.8pt, draw=ThesisNeutral!65]
        (\x,-0.48) -- (\x,-0.22);
      \node[font=\scriptsize\sffamily, text=ThesisNeutral, anchor=north]
        at (\x,-0.50) {\date};
    }
    \node[font=\scriptsize\sffamily\bfseries, text=ThesisNeutral,
      anchor=east] at (-6.95,-0.35) {dossier};

    \node[font=\scriptsize\sffamily\bfseries, text=ThesisPrimaryDark,
      anchor=east] at (-6.95,-1.43) {eCRF};
    \node[field] (f1) at (-5.4,-2.10) {%
      \textbf{date du diagnostic}\\
      12 mai 2018
    };
    \node[field] (f2) at (-1.8,-2.10) {%
      \textbf{traitement, ligne 1}\\
      \textcolor{ThesisPrimaryDark}{\textbf{CHOEP-14}}
    };
    \node[field] (f3) at (1.8,-2.10) {%
      \textbf{réponse, ligne 1}\\
      réponse complète
    };
    \node[field] (f4) at (5.4,-2.10) {%
      \textbf{traitement, ligne 2}\\
      \textcolor{ThesisPrimaryDark}{\textbf{CHOEP-14}}
    };

    \draw[map] (d1.south) -- (f1.north);
    \draw[map, draw=ThesisPrimaryDark] (d2.south) -- (f2.north);
    \draw[map] (d3.south) -- (f3.north);
    \draw[map, draw=ThesisPrimaryDark] (d4.south) -- (f4.north);
  \end{tikzpicture}
  \caption{\textbf{Du dossier longitudinal au formulaire.}
  Dans le dossier synthétique \texttt{lym\_76}, la même chaîne
  \texttt{CHOEP-14} renseigne la première puis la deuxième ligne après
  rechute: le champ dépend de la position de la mention dans la trajectoire,
  pas de sa seule forme locale.}
  \label{fig:p3-task}
\end{figure}
